\documentclass{article}
\usepackage[utf8]{inputenc}
\title{Olá, mundo!}
\author{Bryan Samuel da Silva Cherobini}
\date{16 de Abril de 2026}

\begin{document}

\maketitle

\section{Introdução}


Olá, mundo!
Este documento apresenta a resolução do exercício 1 com o objetivo de revisar a estrutura básica da linguagem LaTeX e a introdução e familiarização do autor com a linguagem.

\section{Desenvolvimento}

Este documento apresenta a resolução do exercício 1 com o objetivo de revisar a estrutura básica da linguagem LaTeX e promover a minha introdução e familiarização com a ferramenta. Por ser o meu primeiro arquivo de autoria própria, ele funciona como um ponto de partida para entender como os comandos funcionam na prática, transformando o código LaTeX em um artigo organizado e estruturado.

O foco principal desse artigo é a adaptação ao novo formato de escrita e linguagem, trocando os editores comuns por uma linguagem de marcação. Além de servir como uma introdução técnica, este trabalho me ajuda a ganhar confiança na criação de documentos originais e garante que eu consiga dominar as bases necessárias para os próximos desafios acadêmicos que utilizarem o LaTeX.

    \subsection{Materiais}
    Para a elaboração deste exercício, utilizei como ferramenta principal o editor [Prism], que permite a escrita e compilação do código em tempo real. Além disso, consultei a documentação básica do LaTeX e modelos acadêmicos para entender a aplicação de comandos de estruturação, garantindo que o documento siga os padrões necessários de formatação técnica.

    \subsection{Métodos}
    O método adotado consistiu na organização do arquivo através de comandos de hierarquia, como seções e subseções, para separar claramente cada parte do trabalho. O processo envolveu a configuração do preâmbulo, a definição de pacotes de linguagem e a escrita estruturada do conteúdo, focando na revisão de elementos que não foram explorados no primeiro exercício, como a organização detalhada do desenvolvimento.

    \subsection{Procedimentos e Ferramentas}

    Para exemplificar a organização de dados em textos científicos, descrevemos abaixo as etapas de um experimento fictício e as ferramentas necessárias para sua execução.

    \noindent \textbf{Etapas do Experimento:}
    \begin{enumerate}
        \item Preparação das amostras em ambiente controlado com temperatura constante.
        \item Aplicação do reagente base e observação das alterações físico-químicas durante trinta minutos.
        \item Coleta dos dados finais e tabulação dos resultados para análise estatística.
    \end{enumerate}

    \noindent \textbf{Ferramentas Utilizadas:}
    \begin{itemize}
        \item Microcomputador com processador de alto desempenho;
        \item Software de edição de texto LaTeX (Overleaf);
        \item Planilhas eletrônicas para organização de dados;
        \item Navegador de internet para pesquisa bibliográfica.
    \end{itemize}

\section{Conclusão}

Pode-se concluir que, este artigo serviu como uma base de conhecimento e introdução a escrita LaTex, após ter escrito todo este documento eu aprendi, Seções, Subtextos, Tabelas(Não foi usado, mas foi revisado), Título, Autores entre outros.

\end{document}